\section{Introduction}

Streaming platforms have revolutionized entertainment, offering convenience and vast content libraries. Among the most popular services, Disney+ is often viewed as a family-friendly platform, while Netflix is known for its diverse content, spanning genres from comedies to mature dramas.

\subsection{Key Questions}

Despite their reputations, both platforms cater to broader audiences. Disney+ features franchises like Marvel and Star Wars, which appeal beyond children, while Netflix includes family-friendly titles. This raises two critical questions:
\begin{itemize}
    \item Does Disney+ predominantly feature content with lower age restrictions?
    \item Are Netflix movies generally of higher quality, based on Rotten Tomatoes scores?
\end{itemize}

\subsection{Relevance of the Analysis}

These questions have practical implications. Families need guidance on selecting platforms with appropriate content, while movie enthusiasts seek high-quality films. Streaming services can also use such insights to refine their content strategies and target audiences more effectively. \cite{Nielsen2021}

\subsection{Approach and Report Structure}

This report uses the dataset of movies from Disney+ and Netflix \cite{dataset}. First, a descriptive analysis explores age ratings and Rotten Tomatoes scores. Then, hypothesis tests evaluate whether observed differences are statistically significant.

The report is organized as follows:
\begin{itemize}
    \item \textbf{Section 2: Problem Description} - Explains the problem and dataset.
    \item \textbf{Section 3: Methods} - Outlines statistical methods and tests.
    \item \textbf{Section 4: Evaluation} - Presents descriptive and inferential analysis results.
    \item \textbf{Section 5: Summary and Discussion} - Concludes with findings and implications.
\end{itemize}
